\documentclass[11pt,a4paper]{article}
\usepackage[margin=2.4cm]{geometry}
\usepackage{amsmath,amssymb}
\usepackage{booktabs}
\usepackage{enumitem}
\usepackage{xcolor}
\usepackage{hyperref}
\usepackage{graphicx}
\usepackage{fancyhdr}
\usepackage{titlesec}
\usepackage{multirow}
\usepackage{siunitx}
\usepackage{bm}

\setlength{\headheight}{14pt}

\pagestyle{fancy}
\fancyhf{}
\rhead{\textit{Methodology}}
\lhead{\textit{Europa Ice Shell Thermal Model}}
\cfoot{\thepage}

\titleformat{\section}{\Large\bfseries}{}{0em}{}
\titleformat{\subsection}{\large\bfseries}{}{0em}{}
\titleformat{\subsubsection}{\normalsize\bfseries\itshape}{}{0em}{}

\newcommand{\dd}{\mathrm{d}}
\newcommand{\pp}{\partial}
\newcommand{\Ra}{\mathrm{Ra}}
\newcommand{\Nu}{\mathrm{Nu}}

\title{\textbf{Model Methodology}\\[0.3em]
\Large Thermal Structure and Uncertainty Quantification\\
of Europa's Ice Shell\\[0.5em]
\large\textit{Full Mathematical Description}}
\author{}
\date{}

\begin{document}
\maketitle
\thispagestyle{fancy}

\tableofcontents
\newpage

%======================================================================
\section{Overview}
%======================================================================

This chapter presents the complete mathematical framework for modelling
the thermal structure of Europa's ice shell. The model comprises five
coupled components:

\begin{enumerate}[leftmargin=2em]
\item A \textbf{1D transient thermal solver} for the radial temperature
      profile through the ice shell (\S\ref{sec:thermal}).
\item A \textbf{rheological model} coupling ice viscosity to temperature,
      grain size, and tidal forcing (\S\ref{sec:rheology}).
\item A \textbf{parameterised stagnant-lid convection scheme} that
      partitions the shell into a conductive lid and convective sublayer
      (\S\ref{sec:convection}).
\item A \textbf{2D latitude-dependent extension} that couples independent
      radial columns through latitude-varying boundary conditions
      (\S\ref{sec:2d}).
\item A \textbf{Monte Carlo uncertainty quantification framework} with
      Bayesian observational constraining (\S\ref{sec:mc}).
\end{enumerate}

%======================================================================
\section{Thermal Solver}
\label{sec:thermal}
%======================================================================

\subsection{Governing Equation}

The model solves the one-dimensional transient heat conduction equation
with volumetric tidal heating:
\begin{equation}
\rho(T)\, c_p(T)\, \frac{\pp T}{\pp t}
  = \frac{1}{\mathcal{G}}\,\frac{\pp}{\pp z}
    \left(\mathcal{G}\, k(T)\, \frac{\pp T}{\pp z}\right)
  + q_{\mathrm{tidal}}(T)
\label{eq:heat}
\end{equation}
where $z$ is depth measured from the surface, $T(z,t)$ is temperature,
and $\mathcal{G}$ is a geometric factor: $\mathcal{G} = 1$ in Cartesian
coordinates (thin shells, $H < \SI{30}{\kilo\metre}$) or
$\mathcal{G} = r^2$ in spherical coordinates (thick shells,
$H \geq \SI{30}{\kilo\metre}$), where $r = R_p - z$ and
$R_p = \SI{1561}{\kilo\metre}$ is Europa's radius.

\subsection{Material Properties}

\subsubsection{Thermal conductivity}

The thermal conductivity of ice~I follows the Howell (2021) parameterisation:
\begin{equation}
k(T) = \frac{567}{T} \quad \si{\watt\per\metre\per\kelvin}
\label{eq:conductivity}
\end{equation}
This inverse dependence on temperature is the single most important
relationship in the model: it produces a factor-of-four conductivity
contrast between the cold polar surface ($T \sim \SI{50}{\kelvin}$,
$k \sim \SI{11}{\watt\per\metre\per\kelvin}$) and the warm equatorial
surface ($T \sim \SI{110}{\kelvin}$,
$k \sim \SI{5}{\watt\per\metre\per\kelvin}$).

Where near-surface porosity is present ($T < T_\phi = \SI{150}{\kelvin}$),
the effective conductivity is reduced:
\begin{equation}
k_{\mathrm{eff}}(T) = k(T)\,(1 - f_{\mathrm{porosity}})
\quad \text{for } T < T_\phi
\label{eq:porosity}
\end{equation}
where $f_{\mathrm{porosity}} \in [0, 0.10]$ is the porosity fraction.

\subsubsection{Density}

Temperature-dependent density follows Howell (2021):
\begin{equation}
\rho(T) = \rho_0\left[1 + \alpha_{th}(T_m - T)\right]
\label{eq:density}
\end{equation}
with reference density $\rho_0 = \SI{917}{\kilo\gram\per\metre\cubed}$
at the melting point $T_m = \SI{273}{\kelvin}$ and thermal expansion
coefficient $\alpha_{th} = \SI{1.6e-4}{\per\kelvin}$.

\subsubsection{Specific heat capacity}

The specific heat capacity follows the temperature-dependent fit of
Giauque \& Stout (1936):
\begin{equation}
c_p(T) = 7.49\,T + 90.0 \quad \si{\joule\per\kilo\gram\per\kelvin}
\label{eq:cp}
\end{equation}
yielding $c_p \approx \SI{839}{\joule\per\kilo\gram\per\kelvin}$ at
$T = \SI{100}{\kelvin}$ and
$c_p \approx \SI{2135}{\joule\per\kilo\gram\per\kelvin}$ at the
melting point.

\subsection{Boundary Conditions}
\label{sec:bc}

\subsubsection{Surface boundary condition}

The surface temperature is prescribed as a Dirichlet condition:
\begin{equation}
T(z=0, t) = T_{\mathrm{surf}}
\label{eq:surf_bc}
\end{equation}
where $T_{\mathrm{surf}}$ is drawn from the Monte Carlo prior
distribution (\S\ref{sec:mc}) or computed from the latitude-dependent
profile (\S\ref{sec:2d}).

Alternatively, a radiative-equilibrium condition may be applied:
\begin{equation}
\varepsilon\,\sigma\,T_s^4
  = (1 - a)\,Q_{\odot}(\phi)
  + k(T_s)\left.\frac{\pp T}{\pp z}\right|_{z=0}
\label{eq:stefan_boltzmann}
\end{equation}
where $\varepsilon = 0.94$ is the surface emissivity,
$\sigma = \SI{5.67e-8}{\watt\per\metre\squared\per\kelvin\tothe{4}}$
the Stefan--Boltzmann constant, $a = 0.64$ the Bond albedo, and
$Q_{\odot}$ the local solar flux. This nonlinear equation is solved
via Newton--Raphson iteration using the linearisation
$T^4 \approx 4T_0^3 T - 3T_0^4$.

\subsubsection{Basal boundary condition: Stefan condition}

At the ice--ocean interface ($z = H$), the temperature is fixed at the
pressure-dependent melting point:
\begin{equation}
T(z=H) = T_m(H) = 273.0 - \SI{7.4e-8}{\kelvin\per\pascal}
  \times \rho\,g\,H
\label{eq:melt}
\end{equation}
where $g = \SI{1.315}{\metre\per\second\squared}$ is Europa's surface
gravity. For a \SI{20}{\kilo\metre} shell, $T_m \approx \SI{272.8}{\kelvin}$.

The shell thickness $H$ evolves according to the Stefan condition
governing the freezing front:
\begin{equation}
\frac{\dd H}{\dd t}
  = \frac{q_{\mathrm{cond}} - q_{\mathrm{ocean}}}{\rho\, L}
\label{eq:stefan}
\end{equation}
where $L = \SI{334000}{\joule\per\kilo\gram}$ is the latent heat of
fusion, $q_{\mathrm{ocean}}$ is the heat flux delivered from the ocean
to the base of the ice, and $q_{\mathrm{cond}}$ is the heat conducted
upward from the base:
\begin{equation}
q_{\mathrm{cond}} = k(T_m)\left.\frac{\pp T}{\pp z}\right|_{z=H}
\label{eq:qcond}
\end{equation}
evaluated using a second-order one-sided finite difference:
\begin{equation}
\left.\frac{\pp T}{\pp z}\right|_{z=H}
  \approx \frac{3T_n - 4T_{n-1} + T_{n-2}}{2\,\Delta z}
\label{eq:grad}
\end{equation}

If $q_{\mathrm{cond}} > q_{\mathrm{ocean}}$, the shell thickens
(freezing); if $q_{\mathrm{cond}} < q_{\mathrm{ocean}}$, it thins
(melting). The model is integrated to thermal equilibrium, defined as
$|\dd H/\dd t| < \SI{e-15}{\metre\per\second}$.

\subsection{Numerical Discretisation}
\label{sec:numerics}

\subsubsection{Spatial discretisation}

The radial domain $[0, H]$ is divided into $n_z$ equally spaced nodes
(default $n_z = 31$ for Monte Carlo production runs). The diffusion
operator is discretised in flux-conservative form using harmonic-mean
averaging of conductivity at half-nodes:
\begin{equation}
k_{i+1/2} = \frac{2\,k_i\,k_{i+1}}{k_i + k_{i+1}}
\label{eq:harmonic}
\end{equation}
This ensures flux conservation across the conductive--convective
interface where $k_{\mathrm{eff}}$ may jump by an order of magnitude
when the Nusselt enhancement is applied (\S\ref{sec:convection}).

\subsubsection{Time integration: Crank--Nicolson with Rannacher startup}

The heat equation is advanced in time using the $\theta$-method:
\begin{equation}
\rho\,c_p\,\frac{T^{n+1}_i - T^n_i}{\Delta t}
  = \theta\,\mathcal{L}[T^{n+1}]
  + (1-\theta)\,\mathcal{L}[T^n]
  + q_{\mathrm{tidal}}
\label{eq:theta}
\end{equation}
where $\mathcal{L}$ denotes the discrete spatial diffusion operator.

A Rannacher startup (Rannacher 1984) is employed: the first four
half-steps use $\theta = 1$ (backward Euler) at $\Delta t/2$ to damp
spurious oscillations arising from stiff initial conditions. All
subsequent steps use $\theta = 0.5$ (Crank--Nicolson) at the full
timestep $\Delta t = \SI{e12}{\second}$ ($\sim$31,700~years), giving
second-order temporal accuracy.

This yields a tridiagonal system at each timestep:
\begin{equation}
a_i\,T^{n+1}_{i-1} + b_i\,T^{n+1}_i + c_i\,T^{n+1}_{i+1} = d_i
\label{eq:tridiag}
\end{equation}
solved via the Thomas algorithm in $O(n_z)$ operations.

\subsubsection{Picard iteration for nonlinearity}

Because $k(T)$, $\rho(T)$, $c_p(T)$, and $q_{\mathrm{tidal}}(T)$ all
depend on temperature, the system is nonlinear. At each timestep, up to
three Picard iterations are performed: material properties are evaluated
at the current temperature estimate, the tridiagonal system is assembled
and solved, and the new temperature is used to update properties.
Convergence is declared when
$\max_i |T^{(m+1)}_i - T^{(m)}_i| < \SI{0.01}{\kelvin}$.

%======================================================================
\section{Rheology and Tidal Dissipation}
\label{sec:rheology}
%======================================================================

\subsection{Composite Diffusion Creep Viscosity}

The viscosity of polycrystalline ice is computed from the composite
diffusion creep law of Howell (2021), which includes both volume
diffusion and grain-boundary diffusion:
\begin{equation}
\eta(T, d) = \frac{1}{2}\left[
  \frac{42\,\Omega}{R\,T\,d^2}
  \left(D_v(T) + \frac{\pi\,\delta}{d}\,D_b(T)\right)
\right]^{-1}
\label{eq:viscosity}
\end{equation}
where $\Omega = \SI{1.97e-5}{\metre\cubed\per\mole}$ is the molar
volume of ice, $R = \SI{8.314}{\joule\per\mole\per\kelvin}$ the gas
constant, $d$ the grain size, and $\delta = \SI{7.13e-10}{\metre}$
the grain-boundary width. The diffusion coefficients follow Arrhenius
laws:
\begin{align}
D_v(T) &= D_{0v}\,\exp\!\left(-\frac{Q_v}{R\,T}\right) \\
D_b(T) &= D_{0b}\,\exp\!\left(-\frac{Q_b}{R\,T}\right)
\label{eq:diffusion}
\end{align}
with parameters listed in Table~\ref{tab:rheology}. Viscosity is
clipped to the range $[10^{12},\, 10^{25}]$~Pa\,s for numerical
stability.

\begin{table}[h]
\centering
\caption{Rheological parameters (Howell 2021).}
\label{tab:rheology}
\begin{tabular}{llrl}
\toprule
\textbf{Parameter} & \textbf{Symbol} & \textbf{Value} & \textbf{Unit} \\
\midrule
Volume diffusion prefactor  & $D_{0v}$ & $9.1 \times 10^{-4}$ & m$^2$/s \\
Boundary diffusion prefactor & $D_{0b}$ & $8.4 \times 10^{-4}$ & m$^2$/s \\
Volume activation energy    & $Q_v$    & 59,400 & J/mol \\
Boundary activation energy  & $Q_b$    & 49,000 & J/mol \\
Molar volume                & $\Omega$ & $1.97 \times 10^{-5}$ & m$^3$/mol \\
Grain-boundary width        & $\delta$ & $7.13 \times 10^{-10}$ & m \\
Ice shear modulus            & $\mu$    & $3.3 \times 10^{9}$ & Pa \\
\bottomrule
\end{tabular}
\end{table}

\subsection{Tidal Heating: Maxwell Rheology}

Under the Maxwell (spring-dashpot) model, volumetric tidal dissipation is:
\begin{equation}
q_{\mathrm{tidal}}^{\mathrm{Max}}(T) =
  \frac{\varepsilon_0^2\,\omega^2\,\eta(T)}
       {2\left[1 + \dfrac{\omega^2\,\eta(T)^2}{\mu^2}\right]}
\label{eq:maxwell}
\end{equation}
where $\varepsilon_0$ is the tidal strain amplitude, $\omega =
\SI{2.047e-5}{\radian\per\second}$ is Europa's orbital frequency,
and $\mu = \SI{3.3e9}{\pascal}$ is the ice shear modulus. The factor
of $1/2$ arises from time-averaging $\sin^2(\omega t)$ over one
orbital period.

Maxwell dissipation peaks sharply at $\omega\tau = 1$ where $\tau =
\eta/\mu$ is the Maxwell relaxation time. Only a narrow temperature
band satisfies this resonance condition; ice that is too cold
($\omega\tau \gg 1$, elastic response) or too warm
($\omega\tau \ll 1$, viscous flow) dissipates negligibly.

\subsection{Tidal Heating: Andrade Rheology}

The Andrade model (McCarthy et al.\ 2011; Renaud \& Henning 2018)
incorporates transient creep from sub-grain dislocation motion and
grain-boundary sliding, producing broadband dissipation absent from
the Maxwell model.

The complex compliance is:
\begin{equation}
J^*(\omega) = \underbrace{\frac{1}{\mu}}_{\text{elastic}}
  + \underbrace{\frac{\Gamma(1+\alpha)}{\mu}\,
    (\omega\tau\zeta)^{-\alpha}\,
    e^{-i\alpha\pi/2}}_{\text{Andrade transient creep}}
  - \underbrace{\frac{i}{\omega\,\eta}}_{\text{viscous}}
\label{eq:andrade_compliance}
\end{equation}
where $\alpha = 0.2$ is the Andrade creep exponent, $\zeta = 1.0$
is the normalisation factor, $\tau = \eta/\mu$ is the relaxation time,
and $\Gamma$ is the gamma function. The real and imaginary parts are:
\begin{align}
J_{\mathrm{re}} &= \frac{1}{\mu} + \frac{\Gamma(1+\alpha)}{\mu}\,
  (\omega\tau\zeta)^{-\alpha}\,\cos\!\left(\frac{\alpha\pi}{2}\right)
\label{eq:Jre} \\
J_{\mathrm{im}} &= \frac{1}{\omega\,\eta} + \frac{\Gamma(1+\alpha)}{\mu}\,
  (\omega\tau\zeta)^{-\alpha}\,\sin\!\left(\frac{\alpha\pi}{2}\right)
\label{eq:Jim}
\end{align}

The dissipative (imaginary) part of the shear modulus is:
\begin{equation}
G'' = \frac{J_{\mathrm{im}}}{J_{\mathrm{re}}^2 + J_{\mathrm{im}}^2}
\label{eq:Gimag}
\end{equation}

Volumetric tidal heating under Andrade rheology is then:
\begin{equation}
q_{\mathrm{tidal}}^{\mathrm{And}}(T) = \tfrac{1}{2}\,\omega\,
  \varepsilon_0^2\,G''(\omega, T)
\label{eq:andrade_heating}
\end{equation}

The Andrade term $(\omega\tau)^{-\alpha}$ decays as a gentle power law
($\alpha = 0.2$), producing significant dissipation across a much wider
range of viscosities (and therefore temperatures) than the sharp Maxwell
peak. This is the primary reason for adopting Andrade rheology: laboratory
measurements on polycrystalline ice confirm broadband dissipation
behaviour inconsistent with the single-timescale Maxwell model.

%======================================================================
\section{Stagnant-Lid Convection Parameterisation}
\label{sec:convection}
%======================================================================

Europa's ice shell is expected to convect in the stagnant-lid regime,
where a rigid conductive lid overlies a slowly overturning convective
interior. Rather than resolving convective flow, the model parameterises
the heat transport enhancement using scaling laws from numerical
experiments (Solomatov \& Moresi 2000; Green et al.\ 2021;
Deschamps \& Vilella 2021).

\subsection{Conductive--Convective Interface Detection}

At each timestep, the model scans the temperature profile from the
surface to find the critical rheological transition temperature $T_c$
below which the viscosity contrast is too large for convective
participation.

\subsubsection{Interior temperature (Deschamps \& Vilella 2021, Eq.~18)}

\begin{equation}
T_i = B\left[\sqrt{1 + \frac{2}{B}(T_m - c_2\,\Delta T_s)} - 1\right]
\label{eq:Ti}
\end{equation}
where $B = Q_v / (2\,R\,c_1)$, $c_1 = 1.43$, $c_2 = -0.03$, and
$\Delta T_s = T_m - T_{\mathrm{surf}}$.

\subsubsection{Viscous temperature scale (Green et al.\ 2021)}

The viscous temperature scale $\Delta T_v$ characterises the temperature
interval over which viscosity changes by one $e$-fold:
\begin{equation}
\Delta T_v = -\frac{\eta(T_i)}{\dd\eta/\dd T\big|_{T_i}}
  = \frac{R\,T_i^2}{A_v\,T_m}
\label{eq:DTv}
\end{equation}
where $A_v = Q_v / (N\,R\,T_m)$ and $N$ is the stress exponent.

\subsubsection{Lid-base temperature}

The conductive lid base temperature is:
\begin{equation}
T_c = T_i - \theta_{\mathrm{lid}}\,\Delta T_v
\label{eq:Tc}
\end{equation}
where $\theta_{\mathrm{lid}} = 2.24$ (Green et al.\ 2021, following
Deschamps \& Vilella 2021).

\subsubsection{Interface depth}

The conductive lid thickness $D_{\mathrm{cond}}$ is obtained by scanning
the numerical temperature profile for the first depth $z_c$ where
$T(z) \geq T_c$, with linear interpolation between grid nodes:
\begin{equation}
z_c = z_{j-1} + \frac{T_c - T_{j-1}}{T_j - T_{j-1}}\,(z_j - z_{j-1})
\label{eq:zc}
\end{equation}

The convective sublayer thickness is $D_{\mathrm{conv}} = H - z_c$.

\subsection{Rayleigh and Nusselt Numbers}

The Rayleigh number of the convective sublayer is:
\begin{equation}
\Ra = \frac{\rho\,g\,\alpha_{th}\,\Delta T_{\mathrm{conv}}\,
  D_{\mathrm{conv}}^3}{\kappa\,\bar{\eta}}
\label{eq:Ra}
\end{equation}
where $\Delta T_{\mathrm{conv}} = T_m - T_c$, $\kappa = k/(\rho\,c_p)$
is the thermal diffusivity, and $\bar{\eta}$ is the viscosity evaluated
at the mean convective temperature. Convection is active when
$\Ra \geq \Ra_{\mathrm{crit}} = 1000$.

For supercritical flow, the Nusselt number is:
\begin{equation}
\Nu = C\,\Ra^{\xi}\left(\frac{T_i - T_c}{\Delta T_{\mathrm{conv}}}\right)^{\!\zeta}
\label{eq:Nu}
\end{equation}
with $C = 0.3446$, $\xi = 1/3$ (Solomatov \& Moresi 2000), and
$\zeta = 4/3$ accounting for internal tidal heating
(Green et al.\ 2021). For subcritical conditions ($\Ra < \Ra_{\mathrm{crit}}$),
$\Nu = 1$ (pure conduction).

\subsection{Effective Conductivity}

Within the convective sublayer, the thermal conductivity is enhanced
by the Nusselt factor:
\begin{equation}
k_{\mathrm{eff}}(z) = \begin{cases}
  k(T(z))            & z < z_c \quad \text{(conductive lid)} \\
  \Nu \cdot k(T(z))  & z \geq z_c \quad \text{(convective sublayer)}
\end{cases}
\label{eq:keff}
\end{equation}

The 2D column-array model applies this same Nusselt enhancement
directly in every latitude column, with no additional surface-temperature
ramp or cold-column suppression factor. This keeps the 1D and 2D
convection criteria identical.

%======================================================================
\section{Latitude-Dependent Extension}
\label{sec:2d}
%======================================================================

The 1D model is extended to a 2D axisymmetric configuration by
coupling $N_\phi$ independent radial columns (default $N_\phi = 19$)
at latitudes $\phi \in [0, \pi/2]$ (equator to pole), each with
latitude-appropriate boundary conditions.

\subsection{Surface Temperature Profile}

The annual-mean surface temperature varies with latitude as:
\begin{equation}
T_s(\phi) = \left[
  (T_{\mathrm{eq}}^4 - T_{\mathrm{floor}}^4)\,\cos^p\!\phi
  + T_{\mathrm{floor}}^4
\right]^{1/4}
\label{eq:Tsurf_lat}
\end{equation}
where $T_{\mathrm{eq}} = \SI{110}{\kelvin}$ is the equatorial
temperature, $T_{\mathrm{floor}} = \SI{46}{\kelvin}$ is the polar
floor (Ojakangas \& Stevenson 1989), and $p = 1.25$ is an exponent
calibrated against the seasonal energy balance model of Ashkenazy (2019),
reducing the RMS temperature error from \SI{3.01}{\kelvin} ($p = 1$) to
\SI{0.70}{\kelvin}.

\subsection{Tidal Strain Latitude Profile}

Tidal strain varies with latitude following the eccentricity-tide
dissipation pattern of Beuthe (2013). The default whole-shell
(mantle--core) pattern gives monotonically increasing strain toward
the poles:
\begin{equation}
\varepsilon_0(\phi)
  = \varepsilon_{\mathrm{eq}}\,\sqrt{1 + c\,\sin^2\!\phi}
\label{eq:strain_lat}
\end{equation}
where
\begin{equation}
c = \left(\frac{\varepsilon_{\mathrm{pole}}}
  {\varepsilon_{\mathrm{eq}}}\right)^{\!2} - 1
\end{equation}
With $\varepsilon_{\mathrm{eq}} = \num{6.0e-6}$ and
$\varepsilon_{\mathrm{pole}} = \num{1.2e-5}$ (a 2:1 pole-to-equator
ratio; Tobie et al.\ 2003), $c = 3$.

Two alternative patterns are also implemented:
\begin{itemize}[leftmargin=2em]
\item \textbf{Shell-dominated} (thin-shell membrane):
  $\varepsilon_0(\phi) = \varepsilon_{\mathrm{pole}}\,
  \sqrt{1 + c'\cos^2\!\phi}$ — monotonically increasing toward the
  equator.
\item \textbf{Non-monotonic} (degree-4 amplification):
  mantle--core $\times [1 + A_m\,\sin^2(2\phi)]$ — peak at
  mid-latitudes ($A_m = 0.3$ default).
\end{itemize}

\subsection{Ocean Heat Transport Scenarios}

The basal heat flux delivered by the ocean to the ice base varies with
latitude under three competing scenarios, parameterised following
Lemasquerier et al.\ (2023):

\begin{enumerate}[leftmargin=2em]
\item \textbf{Uniform transport:}
\begin{equation}
q_{\mathrm{ocean}}(\phi) = q_{\mathrm{mean}}
\end{equation}

\item \textbf{Polar-enhanced} (Lemasquerier et al.\ 2023):
\begin{equation}
q_{\mathrm{ocean}}(\phi) = q_{\mathrm{mean}}\,
  \frac{1 + a\,\sin^2\!\phi}{1 + a/3}
\label{eq:q_polar}
\end{equation}

\item \textbf{Equator-enhanced} (Soderlund et al.\ 2014):
\begin{equation}
q_{\mathrm{ocean}}(\phi) = q_{\mathrm{mean}}\,
  \frac{1 + a\,\cos^2\!\phi}{1 + 2a/3}
\label{eq:q_equator}
\end{equation}
\end{enumerate}

The amplitude parameter $a$ is derived from the mantle tidal fraction
$q^*$ (Lemasquerier et al.\ 2023), which measures the fraction of basal
heating attributable to tidal dissipation in the silicate mantle
($q^* \in [0, 0.91]$):
\begin{equation}
a = \frac{3\,q^*}{3 - n\,q^*}
\end{equation}
where $n = 1$ for polar-enhanced and $n = 2$ for equator-enhanced
patterns. The normalisation denominators in Eqs.~\ref{eq:q_polar}--\ref{eq:q_equator}
ensure that the area-weighted mean heat flux is exactly $q_{\mathrm{mean}}$.

\subsection{Lateral Heat Diffusion}

Explicit lateral diffusion between adjacent columns is applied in
geographic coordinates:
\begin{equation}
\frac{\pp T}{\pp t}\bigg|_{\mathrm{lat}}
  = \frac{k}{R_p^2\,\cos\phi}\,\frac{\pp}{\pp\phi}
    \left(\cos\phi\,\frac{\pp T}{\pp\phi}\right)
\label{eq:lateral}
\end{equation}

At the equator ($\phi = 0$), a symmetry condition
$\pp T/\pp\phi = 0$ is applied. At the pole ($\phi = \pi/2$), where
$\cos\phi \to 0$, L'H\^{o}pital's rule reduces the operator to:
\begin{equation}
\frac{\pp T}{\pp t}\bigg|_{\mathrm{pole}}
  = \frac{2\,k}{R_p^2}\,\frac{\pp^2 T}{\pp\phi^2}
\end{equation}

The thermal diffusion timescale across Europa's radius is
$\tau \sim R_p^2/\kappa \sim 200$~Gyr, far exceeding the
model's integration window of ${\sim}16$~Myr. Consequently,
lateral conduction is negligible; the latitude structure is determined
entirely by the latitude-varying boundary conditions and tidal forcing.

%======================================================================
\section{Monte Carlo Uncertainty Quantification}
\label{sec:mc}
%======================================================================

\subsection{Forward Propagation}

For each Monte Carlo realisation, the model:
\begin{enumerate}[leftmargin=2em]
\item Draws parameter values from the prior distributions
      (Table~\ref{tab:priors}).
\item Initialises a 1D ice column with a linear temperature profile.
\item Integrates Eq.~\ref{eq:heat} to thermal equilibrium using the
      Stefan condition (Eq.~\ref{eq:stefan}) to evolve shell thickness.
\item Records the equilibrium conductive lid thickness $D_{\mathrm{cond}}$,
      convective sublayer thickness $D_{\mathrm{conv}}$, total shell
      thickness $H$, Rayleigh and Nusselt numbers, and all sampled
      parameter values.
\end{enumerate}

Production ensembles use $N = 5{,}000$--$15{,}000$ realisations,
parallelised across CPU cores. Physical filters reject realisations with
$H < \SI{0.5}{\kilo\metre}$ or $H > 0.99\,D_{\mathrm{H_2O}}$ (water
layer thickness).

\subsection{Prior Distributions}
\label{sec:priors}

Table~\ref{tab:priors} lists the prior distributions for all sampled
parameters, based on the 2026 parameter audit
(PARAMETER\_PRIOR\_AUDIT\_2026).

\begin{table}[h]
\centering
\caption{Monte Carlo prior distributions (2026 audited).}
\label{tab:priors}
\small
\begin{tabular}{llll}
\toprule
\textbf{Parameter} & \textbf{Distribution} & \textbf{Range / Clip} & \textbf{Sensitivity} \\
\midrule
$q_{\mathrm{basal}}$ & Uniform & [5, 25] mW/m$^2$ & Medium \\
$d_{\mathrm{grain}}$ & LogNormal($\SI{0.6}{\milli\metre}$, $\sigma = 0.35$~dex)
  & [0.05, 3.0] mm & \textbf{High} \\
$\varepsilon_0$ & LogNormal($1.2{\times}10^{-5}$, $\sigma = 0.3$~dex)
  & [$2{\times}10^{-6}$, $3.4{\times}10^{-5}$] & High \\
$\mu_{\mathrm{ice}}$ & Normal($3.5 \pm 0.5$ GPa) & [2, 5] GPa & High \\
$T_{\mathrm{surf}}$ & Normal($104 \pm 7$ K) & [80, 120] K & Medium \\
$D_{\mathrm{H_2O}}$ & Normal($127 \pm 21$ km) & [80, 200] km & Low \\
$Q_v$ & Normal(59.4 kJ/mol, 5\%) & --- & \textbf{Dominant} \\
$Q_b$ & Normal(49.0 kJ/mol, 5\%) & --- & Low \\
$D_{0v}$ & Normal($9.1{\times}10^{-4}$, 3.3\%) & --- & Low \\
$D_{0b}$ & Normal($8.4{\times}10^{-4}$, 3.3\%) & --- & Low \\
$H_{\mathrm{rad}}$ & $|\mathcal{N}(4.5, 1.0)|$ pW/kg & $> 0$ & Low \\
$f_{\mathrm{porosity}}$ & Uniform & [0, 0.10] & Negligible \\
$f_{\mathrm{salt}}$ & Fixed at 0 & --- & --- \\
\bottomrule
\end{tabular}
\end{table}

\subsubsection{Equatorial overrides}

For equatorial-proxy runs, $T_{\mathrm{surf}}$ is replaced by
$\mathcal{N}(\SI{110}{\kelvin}, \SI{5}{\kelvin})$ clipped to
$[\SI{95}{\kelvin}, \SI{120}{\kelvin}]$, and tidal strain by
LogNormal($\num{6e-6}$, $\sigma = 0.2$~dex) clipped to
$[\num{2e-6}, \num{2e-5}]$. An ocean heat scaling factor
$f_{\mathrm{ocean}} \in \{0.55, 0.67, 1.0, 1.2, 1.5\}$ multiplies
the tidal component of $q_{\mathrm{basal}}$.

\subsubsection{Polar overrides}

For polar-proxy runs, $T_{\mathrm{surf}}$ is
$\mathcal{N}(\SI{50}{\kelvin}, \SI{5}{\kelvin})$ clipped to
$[\SI{35}{\kelvin}, \SI{70}{\kelvin}]$, and tidal strain is
LogNormal($\num{1.2e-5}$, $\sigma = 0.2$~dex) clipped to
$[\num{5e-6}, \num{5e-5}]$.

\subsection{Sobol Sensitivity Analysis}

Global sensitivity is quantified via Sobol indices (Sobol' 1993),
computed using Saltelli's extension of the Monte Carlo estimator with
$N = 512$ base samples. The first-order index $S_i$ measures the
fraction of output variance attributable to parameter $i$ alone; the
total-order index $S_{T_i}$ includes all interactions involving
parameter $i$.

\subsection{Bayesian Observational Constraining}
\label{sec:bayes}

The Juno Microwave Radiometer (MWR) provides a constraint on the
conductive lid thickness at mid-latitudes:
$D_{\mathrm{cond}}^{\mathrm{obs}} = 29 \pm \SI{10}{\kilo\metre}$
(Levin et al.\ 2026). Rather than re-running the forward model, we
apply importance reweighting to the existing Monte Carlo ensemble.

Each realisation $i$ receives a weight:
\begin{equation}
w_i = \frac{1}{\sqrt{2\pi}\,\sigma_{\mathrm{tot}}}\,
  \exp\!\left[-\frac{1}{2}
  \left(\frac{D_{\mathrm{cond},i} - D_{\mathrm{cond}}^{\mathrm{obs}}}
       {\sigma_{\mathrm{tot}}}\right)^{\!2}\right]
\label{eq:weights}
\end{equation}
where $\sigma_{\mathrm{tot}} = \sqrt{\sigma_{\mathrm{obs}}^2 +
\sigma_{\mathrm{model}}^2}$ combines the observational uncertainty
($\sigma_{\mathrm{obs}} = \SI{10}{\kilo\metre}$) and a model
discrepancy term ($\sigma_{\mathrm{model}} = \SI{3}{\kilo\metre}$).

Weighted posterior percentiles are computed via the sorted cumulative
weight method:
\begin{equation}
\hat{q}_p = x_{(\pi)}
\quad\text{where}\quad
\sum_{j \leq \pi} w_{(j)} = p \cdot \sum_j w_j
\end{equation}

The effective sample size (ESS) monitors the quality of the
reweighting:
\begin{equation}
\mathrm{ESS} = \frac{\left(\sum_i w_i\right)^2}{\sum_i w_i^2}
\label{eq:ess}
\end{equation}
An iterative scheme tightens the priors on $q_{\mathrm{basal}}$ and
$d_{\mathrm{grain}}$ using posterior percentiles, re-runs the ensemble
with constrained priors, and iterates until convergence (posterior
median within the target range and ESS fraction $> 10\%$).

Bayes factors for model comparison between ocean transport scenarios
$\mathcal{M}_j$ are computed as the ratio of marginal likelihoods:
\begin{equation}
\mathrm{BF}_{jk}
  = \frac{P(D \mid \mathcal{M}_j)}{P(D \mid \mathcal{M}_k)}
  = \frac{\frac{1}{N_j}\sum_{i=1}^{N_j} w_i^{(j)}}
         {\frac{1}{N_k}\sum_{i=1}^{N_k} w_i^{(k)}}
\label{eq:bayes_factor}
\end{equation}

%======================================================================
\section{Summary of Key Physical Constants}
\label{sec:constants}
%======================================================================

\begin{table}[h]
\centering
\caption{Planetary and thermal constants.}
\label{tab:constants}
\begin{tabular}{llrl}
\toprule
\textbf{Parameter} & \textbf{Symbol} & \textbf{Value} & \textbf{Unit} \\
\midrule
Europa radius         & $R_p$           & 1,561,000  & m \\
Surface gravity       & $g$             & 1.315      & m/s$^2$ \\
Orbital frequency     & $\omega$        & $2.047 \times 10^{-5}$ & rad/s \\
Eccentricity          & $e$             & 0.0101     & --- \\
Melting temperature   & $T_m$           & 273        & K \\
Latent heat           & $L$             & 334,000    & J/kg \\
Reference density     & $\rho_0$        & 917        & kg/m$^3$ \\
Thermal expansion     & $\alpha_{th}$   & $1.6 \times 10^{-4}$ & K$^{-1}$ \\
Clausius--Clapeyron   & $\dd T_m/\dd P$ & $-7.4 \times 10^{-8}$ & K/Pa \\
Critical Rayleigh no. & $\Ra_{\mathrm{crit}}$ & 1000  & --- \\
Nu prefactor          & $C$             & 0.3446     & --- \\
Nu exponent           & $\xi$           & $1/3$      & --- \\
Internal heating exp. & $\zeta$         & $4/3$      & --- \\
Lid ratio             & $\theta_{\mathrm{lid}}$ & 2.24 & --- \\
Andrade exponent      & $\alpha$        & 0.2        & --- \\
\bottomrule
\end{tabular}
\end{table}

%======================================================================
\section{Computational Implementation}
\label{sec:implementation}
%======================================================================

The model is implemented in Python using NumPy for array operations and
SciPy for the tridiagonal solver (\texttt{scipy.linalg.solve\_banded}).
Monte Carlo realisations are parallelised via Python's
\texttt{multiprocessing} module. Default numerical parameters are listed
in Table~\ref{tab:numerical}.

\begin{table}[h]
\centering
\caption{Default numerical parameters.}
\label{tab:numerical}
\begin{tabular}{lrl}
\toprule
\textbf{Parameter} & \textbf{Value} & \textbf{Description} \\
\midrule
$n_z$               & 31        & Radial nodes per column \\
$N_\phi$            & 19        & Latitude columns (2D mode) \\
$\Delta t$          & $10^{12}$~s & Timestep ($\sim$31,700~yr) \\
$t_{\mathrm{max}}$  & $10^{15}$~s & Max integration time ($\sim$31.7~Myr) \\
Rannacher steps     & 4         & BE half-steps at startup \\
Picard iterations   & 3         & Max per timestep \\
Equilibrium tol.    & $10^{-15}$~m/s & $|\dd H/\dd t|$ halt criterion \\
MC samples          & 5,000--15,000 & Per scenario \\
\bottomrule
\end{tabular}
\end{table}

\end{document}
